% GENERATED FILE. Edit canonical lessons, then run npm run generate:books.
% canonical-source-hash: fnv1a64:427149822bd6ec50
% canonical-lessons: HI-C148-kabutar, HI-C148-tota, HI-C148-ullu, HI-C148-titli, HI-C148-makkhi

\chapter{Pigeon, Parrot, Owl, Butterfly, Fly}
\label{ch:hi-pigeon-parrot-owl-butterfly-fly}
\hlchaptermodality{148}

\begin{chapteropening}
\textbf{By the end of this chapter:} \emph{I can say a pigeon, a parrot, an owl, a butterfly and a fly.}

Complete the last lesson of chapter 148: i can say a pigeon, a parrot, an owl, a butterfly and a fly.
\end{chapteropening}

\section[kabūtar]{\hi{कबूतर} (kabūtar) — a pigeon}

\label{lesson:HI-C148-kabutar}



\begin{quote}
Before the new one: say the Hindi for a wolf, then the Hindi for a crow.
\end{quote}

\subsection*{You'll want to know: \hi{कबूतर}}
\textbf{\hi{कबूतर}} — \emph{kabūtar} — \textquotedblleft{}a pigeon\textquotedblright{}.

\begin{grammarlens}[title={What you've built}]
One more word for this chapter.
\end{grammarlens}

\subsection*{Your turn}
\begin{itemize}
\raggedright
  \item \emph{Say it:} \emph{kabūtar}
  \item \emph{Say it:} \emph{kabūtar}, once more
  \item \emph{Say it:} say \emph{kabūtar}
  \item \emph{Recall:} say the Hindi for a wolf, then the Hindi for a crow, then say \emph{kabūtar} again
\end{itemize}

\begin{tcolorbox}[breakable,colback=teal!4,colframe=teal!35!black,title={Before you move on}]
What does \hi{कबूतर} mean? (\textquotedblleft{}A pigeon\textquotedblright{}.) Say it once more.
\end{tcolorbox}
\section[totā]{\hi{तोता} (totā) — a parrot}

\label{lesson:HI-C148-tota}



\begin{quote}
Before the new one: say the Hindi for a crow, then the Hindi for a pigeon.

Order three things two ways: \textbf{\hi{पहली}}, \textbf{\hi{दूसरी}}, \textbf{\hi{तीसरी}}, then \textbf{\hi{एक} … \hi{दूसरा}}. Count on to \textbf{\hi{नौवाँ}} and \textbf{\hi{दसवाँ}}.
\end{quote}

\subsection*{You'll want to know: \hi{तोता}}
\textbf{\hi{तोता}} — \emph{totā} — \textquotedblleft{}a parrot\textquotedblright{}.

\begin{grammarlens}[title={What you've built}]
One more word for this chapter.
\end{grammarlens}

\subsection*{Your turn}
\begin{itemize}
\raggedright
  \item \emph{Say it:} \emph{totā}
  \item \emph{Say it:} \emph{totā}, once more
  \item \emph{Say it:} say \emph{totā}
  \item \emph{Recall:} say the Hindi for a crow, then the Hindi for a pigeon, then say \emph{totā} again
\end{itemize}

\begin{tcolorbox}[breakable,colback=teal!4,colframe=teal!35!black,title={Before you move on}]
What does \hi{तोता} mean? (\textquotedblleft{}A parrot\textquotedblright{}.) Say it once more.
\end{tcolorbox}
\section[ullū]{\hi{उल्लू} (ullū) — an owl}

\label{lesson:HI-C148-ullu}



\begin{quote}
Before the new one: say the Hindi for a pigeon, then the Hindi for a parrot.
\end{quote}

\subsection*{You'll want to know: \hi{उल्लू}}
\textbf{\hi{उल्लू}} — \emph{ullū} — \textquotedblleft{}an owl\textquotedblright{}.

\begin{grammarlens}[title={What you've built}]
One more word for this chapter.
\end{grammarlens}

\subsection*{Your turn}
\begin{itemize}
\raggedright
  \item \emph{Say it:} \emph{ullū}
  \item \emph{Say it:} \emph{ullū}, once more
  \item \emph{Say it:} say \emph{ullū}
  \item \emph{Recall:} say the Hindi for a pigeon, then the Hindi for a parrot, then say \emph{ullū} again
\end{itemize}

\begin{tcolorbox}[breakable,colback=teal!4,colframe=teal!35!black,title={Before you move on}]
What does \hi{उल्लू} mean? (\textquotedblleft{}An owl\textquotedblright{}.) Say it once more.
\end{tcolorbox}
\section[titlī]{\hi{तितली} (titlī) — a butterfly}

\label{lesson:HI-C148-titli}



\begin{quote}
Before the new one: say the Hindi for a parrot, then the Hindi for an owl.
\end{quote}

\subsection*{You'll want to know: \hi{तितली}}
\textbf{\hi{तितली}} — \emph{titlī} — \textquotedblleft{}a butterfly\textquotedblright{}.

\begin{grammarlens}[title={What you've built}]
One more word for this chapter.
\end{grammarlens}

\subsection*{Your turn}
\begin{itemize}
\raggedright
  \item \emph{Say it:} \emph{titlī}
  \item \emph{Say it:} \emph{titlī}, once more
  \item \emph{Say it:} say \emph{titlī}
  \item \emph{Recall:} say the Hindi for a parrot, then the Hindi for an owl, then say \emph{titlī} again
\end{itemize}

\begin{tcolorbox}[breakable,colback=teal!4,colframe=teal!35!black,title={Before you move on}]
What does \hi{तितली} mean? (\textquotedblleft{}A butterfly\textquotedblright{}.) Say it once more.
\end{tcolorbox}
\section[makkhī]{\hi{मक्खी} (makkhī) — a fly}

\label{lesson:HI-C148-makkhi}



\begin{quote}
Before the new one: say the Hindi for an owl, then the Hindi for a butterfly.
\end{quote}

\subsection*{You'll want to know: \hi{मक्खी}}
\textbf{\hi{मक्खी}} — \emph{makkhī} — \textquotedblleft{}a fly\textquotedblright{}.

\begin{grammarlens}[title={What you've built}]
One more word for this chapter.
\end{grammarlens}

\subsection*{Your turn}
\begin{itemize}
\raggedright
  \item \emph{Say it:} \emph{makkhī}
  \item \emph{Say it:} \emph{makkhī}, once more
  \item \emph{Say it:} say \emph{makkhī}
  \item \emph{Recall:} say the Hindi for an owl, then the Hindi for a butterfly, then say \emph{makkhī} again
\end{itemize}

\begin{tcolorbox}[breakable,colback=teal!4,colframe=teal!35!black,title={Before you move on}]
What does \hi{मक्खी} mean? (\textquotedblleft{}A fly\textquotedblright{}.) Say it once more.
\end{tcolorbox}
